\documentclass[11pt]{article}

\usepackage[a4paper,margin=2.5cm]{geometry}
\usepackage{amsmath,amssymb}
\usepackage{graphicx}
\usepackage{booktabs}
\usepackage{titling}
\usepackage{titlesec}
\usepackage{hyperref}
\usepackage{xcolor}
\usepackage{listings}
\usepackage{slashed}

\setlength{\droptitle}{-3em}
\pretitle{\begin{center}\LARGE\bfseries}
\posttitle{\par\end{center}\vskip 0.4em}
\preauthor{\begin{center}\large}
\postauthor{\par\end{center}}
\predate{\begin{center}\small}
\postdate{\par\end{center}}

% Section titles: "<num>. <TITLE IN UPPERCASE>", no bold, centred
\titleformat{\section}[block]
  {\normalfont\normalsize\centering}{\thesection.}{0.5em}{\MakeUppercase}
\titleformat{\subsection}[block]
  {\normalfont\normalsize\itshape\centering}{\thesubsection.}{0.5em}{}
\titleformat{name=\section,numberless}[block]
  {\normalfont\normalsize\centering}{}{0pt}{\MakeUppercase}

\hypersetup{colorlinks=true,linkcolor=blue!50!black,citecolor=blue!50!black,urlcolor=blue!50!black}

\setlength{\emergencystretch}{3em}

\lstset{
  basicstyle=\ttfamily\footnotesize,
  breaklines=true,
  columns=fullflexible,
  showstringspaces=false,
  language=Python,
  keywordstyle=\color{blue!60!black},
  commentstyle=\color{gray},
}

\newcommand{\Dslash}{\slashed{D}}
\newcommand{\NEFF}{\texttt{.neff}}
\newcommand{\PyT}{\textsc{PyTorch}}
\newcommand{\NeuronX}{\textsc{Neuron}}

\title{LQCD-Neuron: Compiling Lattice QCD Kernels\\
       to AWS Trainium and Inferentia via PyTorch--XLA}

\author{Jo\~ao Galego\\
\small Independent Researcher\\
\small \texttt{jgalego1990@gmail.com}}

\date{\today}

\begin{document}

\maketitle

\begin{center}
  \begin{minipage}{0.85\textwidth}
    \begin{center}\bfseries Abstract\end{center}
    \small
We present \textsc{LQCD-Neuron}, a Python library that runs Wilson- and
Clover-Wilson Lattice QCD on AWS Trainium (Trn1) and Inferentia (Inf2)
accelerators. The hardware exposes no general-purpose compute API
(no CUDA, HIP, or OpenCL); the only programming model is to compile a
PyTorch / XLA graph ahead of time with \texttt{neuronx-cc} and execute
the resulting NeuronCore Executable File Format (\NEFF) binary. We
re-express every hot-path lattice kernel - the Wilson Dslash, the
clover term, BLAS-1 reductions, and CG/BiCGStab solvers - as a pure
\texttt{torch.nn.Module}, and rely on three compile-time
transformations (gauge baking, fused $12{\times}12$ spin--colour
matvecs, multi-RHS batching) plus a host-side data-parallel shard
across all on-chip NeuronCores. On a single-core Inf2 NeuronCore-v2
in BF16 we obtain up to $13.0\times$ throughput over a tuned FP32 CPU
baseline for the Wilson Dslash on $V=4{\times}4{\times}4{\times}4$ (two NeuronCores,
$B{=}8$ RHS each), and establish that throughput falls below the CPU
baseline beyond $V=16{\times}8{\times}8{\times}8$, all with no custom kernel code. Source code is available at
\href{https://github.com/JGalego/lqcd-neuron}{github.com/JGalego/lqcd-neuron}
under the Apache 2.0 license.
  \end{minipage}
\end{center}
\vspace{1.5em}


\section{Introduction}
\label{sec:intro}

The cost of a Lattice QCD simulation on heterogeneous hardware is
dominated by repeated applications of the discretised Dirac operator
$D$ to a quark spinor field $\psi$. In the Wilson
formulation~\cite{Wilson:1974sk},
\begin{equation}
(D_W \psi)_x = (4+m)\,\psi_x - \tfrac{1}{2}\sum_{\mu=0}^{3}\!\Big[
  (I-\gamma_\mu)\,U(x,\mu)\,\psi_{x+\hat\mu}
  + (I+\gamma_\mu)\,U^{\dagger}(x{-}\hat\mu,\mu)\,\psi_{x-\hat\mu}\Big],
\label{eq:wilson}
\end{equation}
with gauge links $U(x,\mu)\!\in\!\mathrm{SU}(3)$ and Dirac matrices
$\gamma_\mu$ acting on the four-component spinor. A single application
performs $\mathcal{O}(V)$ tiny dense $3{\times}3$ and $4{\times}4$
contractions, and a Krylov inversion invokes it tens of thousands of
times. On NVIDIA GPUs the standard tool is
\textsc{QUDA}~\cite{Clark:2009wm,Babich:2011np}: hand-written CUDA
kernels, dynamic shapes, runtime auto-tuning, and a multigrid
preconditioner~\cite{Brower:2018kgy}. We ask whether comparable
performance is reachable on AWS-designed ML accelerators --- Trainium
(Trn1) and Inferentia (Inf2) --- which expose a fundamentally
different programming model.

Each NeuronCore-v2 is a purpose-built matrix engine. It offers no
general-purpose compute API, no user-controlled shared memory, and no
runtime kernel launch. Code reaches the hardware along a single path:
a PyTorch (or JAX) computation graph is lowered ahead of time to
XLA HLO~\cite{xla}, the \texttt{neuronx-cc} compiler emits a
\NEFF\ binary, and the Neuron runtime executes it on the NeuronCores.
Tensor shapes must be static at compile time, and the auto-tuning and
runtime kernel re-launch on which QUDA relies are simply not available.
Trainium has been shown to scale deep-learning
training~\cite{aws-trainium}, but to our knowledge has not previously
been targeted for first-principles Lattice QCD.

The central question of this work is therefore: \emph{is the
AoT-compiled, static-graph, matrix-engine programming model
expressive enough for production Lattice QCD primitives, and
competitive in throughput once the graph is shaped to fit the
hardware?} We give an affirmative answer for the Wilson and
Clover-Wilson Dslash and the iterative solvers built on them, with
four composable graph-level optimisations that bring throughput from
below CPU at modest volumes to up to $13.0\times$ above it on a
two-core Inf2 configuration.

\section{Design and architecture}
\label{sec:arch}

\paragraph*{Pure-tensor kernels.}
Every primitive in \textsc{LQCD-Neuron} is a standard
\texttt{torch.nn.Module} operating on dense complex tensors with a
$(T,Z,Y,X,\dots)$ storage order. Gauge links live in
$U\in\mathbb{C}^{T\times Z\times Y\times X\times 4\times N_c\times N_c}$
and spinors in $\psi\in\mathbb{C}^{T\times Z\times Y\times X\times N_s\times N_c}$,
with $N_c=3,\ N_s=4$. Lattice shifts $\psi_{x\pm\hat\mu}$ are
realised by \texttt{torch.roll} with the appropriate axis index, which
the Neuron compiler lowers to in-place data movement on chip. Since
\texttt{neuronx-cc} forbids data-dependent control flow, the
Python loop over the four directions in
Eq.~\eqref{eq:wilson} is statically unrolled at trace time, producing
a pure DAG that the compiler can schedule.

\paragraph*{Gamma matrices.}
The Euclidean Dirac matrices are stored in the DeGrand--Rossi (chiral)
basis~\cite{DeGrandRossi} and registered as non-trainable
\texttt{nn.Module} buffers, so they migrate with the module to
NeuronCore memory under \texttt{torch\_neuronx.trace} and remain
visible to the tracer.

\paragraph*{Host loop / device kernel split.}
Iterative solvers need a convergence test on a scalar reduction at
every iteration, but \texttt{neuronx-cc} only supports static graphs.
We adopt the standard deep-learning split: the Krylov loop runs on the
host (CPU) and only the matvec $D\psi$ is compiled to NeuronCore.
For Conjugate Gradient~\cite{HestenesStiefel}, schematically:
\begin{lstlisting}
for k in range(maxiter):
    Ap    = D_neuron(p, U)        # NeuronCore
    alpha = rho / inner(p, Ap)    # host
    x    += alpha * p
    r    -= alpha * Ap
    if norm(r) < tol: break
\end{lstlisting}
The host overhead per iteration is $\mathcal{O}(V)$ and is dwarfed by
the matvec for any non-trivial lattice; we verified this empirically
against an all-on-host reference.

\paragraph*{CPU fallback.}
Because every kernel is a plain \PyT\ module, the same code runs
identically on a laptop without Neuron hardware. The full unit-test
suite executes on CPU, allowing development and CI without an Inf2
instance.

\section{Compile-time optimisations}
\label{sec:opt}

A naive trace of Eq.~\eqref{eq:wilson} is dominated by overheads that
are absent on a CPU baseline:
(i) the gauge field, $\sim\!3{\times}$ the spinor in size, crosses
PCIe on every call;
(ii) the per-direction kernel is two tiny einsums (a $3{\times}3$
colour multiply and a $4{\times}4$ spin multiply), neither of which
saturates the NeuronCore tensor engine;
(iii) eight \texttt{torch.roll}s per call (four for $\psi$, four for
$U^{\dagger}(x{-}\hat\mu)$) are pure data movement; and
(iv) every NeuronCore call pays a fixed dispatch cost
($\sim\!1\,$ms), independent of problem size. We address these in
four composable steps.

\subsection{Gauge baking}

When a gauge configuration $U$ is supplied at compile time, we register
it as an \texttt{nn.Module} buffer of the traced module rather than as
a graph input. After \texttt{torch\_neuronx.trace}, the resulting
\NEFF\ contains $U$ in NeuronCore-resident memory. At inference time,
only $\psi$ crosses PCIe; the wrapper accepts a $U$ argument for API
compatibility but ignores it. The trade-off is that the \NEFF\ is
specific to a particular $U$, so it is not reused across different
gauge configurations - the right choice for a single inversion or a
chain of solves on one configuration, the wrong choice for HMC.

\subsection{Fused spin--colour kernels}

With $U$ baked in, the per-site, per-direction operator
\begin{equation}
K^{\mathrm{fwd}}_{\mu,x} = (I{-}\gamma_\mu)\!\otimes\! U(x,\mu),\ 
K^{\mathrm{bwd}}_{\mu,x} = (I{+}\gamma_\mu)\!\otimes\! U^{\dagger}(x{-}\hat\mu,\mu)
\end{equation}
can be precomputed once at trace time. We materialise
$K^{\mathrm{fwd/bwd}}_{\mu,x}\in\mathbb{C}^{N_sN_c\times N_sN_c}$ as four
buffers of shape $(4,T,Z,Y,X,12,12)$ stored as real/imaginary float
pairs. Each direction-side then collapses to one $12{\times}12$ matvec
plus a single roll of $\psi$:
\begin{lstlisting}
psi_shifted = torch.roll(psi, sign, dim=mu)
result     -= 0.5 * (K[mu] @ psi_shifted)
\end{lstlisting}
The $12{\times}12$ shape is markedly better suited to the NeuronCore
matrix multiplication unit (MXU) than the original $3{\times}3$ /
$4{\times}4$ einsums. The four backward rolls are absorbed into
$K^{\mathrm{bwd}}$, halving the lattice-shift count. The fused kernel
is bit-equivalent to the reference \texttt{forward()} up to FP32
round-off (relative error $\lesssim 10^{-7}$), and continues to satisfy
the existing $\gamma_5$-Hermiticity, adjoint, and normal-operator
positivity tests.

\subsection{Multi-right-hand-side batching}

Production multi-source propagator inverters already solve $D\,x = b_i$
for $N_s\!\cdot\!N_c=12$ right-hand sides per source, one per
spin\,$\times$\,colour combination. We expose this directly:
\texttt{compile\_dslash\_batched($D$, $\mathrm{shape}$, batch\_size=$B$,
gauge\_field=$U$)} traces the fused operator on
$\psi\in\mathbb{C}^{B\times T\times Z\times Y\times X\times N_s\times N_c}$.
The baked $K$ buffers carry a singleton batch axis that broadcasts
across all $B$ RHS, and the lattice rolls use negative-dim indexing so
the leading batch axis is unaffected. Each direction-side becomes an
$N_b{\times}12{\times}12$ matmul - a much better fit for the MXU's
systolic schedule - and the fixed per-call dispatch cost is
amortised over $B$ problems.

\subsection{Multi-core data-parallel sharding}

A single NeuronCore-v2 is one of $N_{\mathrm{core}}$ cores per
accelerator chip; \texttt{inf2.xlarge} exposes 2,
\texttt{trn1.32xlarge} 32. We replicate the compiled \NEFF\ across
all detected cores via \texttt{torch\_neuronx.DataParallel} and
split a global batch of $N_{\mathrm{core}}\!\cdot\! B_{\mathrm{core}}$
RHS along dim~0:
\begin{lstlisting}
D_mc = compiler.compile_dslash_multicore(
    D, shape, gauge_field=U,
    num_cores=N, per_core_batch_size=B,
)
out  = D_mc(psi_global)   # B_g = N * B
\end{lstlisting}
This composes with the previous three optimisations rather than
replacing them: gauge baking and fused kernels apply per core, the
intra-core multi-RHS batch fills each core's tensor engine, and the
data-parallel shard multiplies by $N_{\mathrm{core}}$.

\section{Results}
\label{sec:results}

\begin{table}[t]
\caption{Wilson Dslash sustained throughput on an \texttt{inf2.8xlarge}
(BF16, gauge baked) versus a tuned CPU baseline (FP32), in Dslash
applications per second over 50 iterations after a 5-iteration warm-up.
Neuron = 1 NeuronCore, 1 RHS; Batched = 1 NeuronCore, $B{=}8$ RHS;
Multicore = 2 NeuronCores, $B{=}8$ RHS each.  Speedup is the
best Neuron column relative to CPU; entries below $1\times$ indicate
the optimised Neuron path has not yet overcome the dispatch and
bandwidth bottleneck at that volume.}
\label{tab:bench}
\centering
\footnotesize
\setlength{\tabcolsep}{4pt}
\begin{tabular}{lrrrrr}
\toprule
Lattice & CPU & Neuron & Batched & Multicore & Speedup \\
\midrule
$4{\times}4{\times}4{\times}4$           &  804.6 & 1358.0 &  8521.5 & \textbf{10431.3} & $13.0\times$ \\
$8{\times}4{\times}4{\times}4$           &  769.5 &  785.1 &  4786.3 &  \textbf{7519.1} &  $9.8\times$ \\
$8{\times}8{\times}4{\times}4$           &  518.1 &  370.1 &  2511.1 &  \textbf{4229.5} &  $8.2\times$ \\
$8{\times}8{\times}8{\times}4$           &  384.0 &  207.6 &  1390.1 &  \textbf{2474.2} &  $6.4\times$ \\
$16{\times}8{\times}8{\times}8$          &  200.9 &   39.3 &   253.7 &    \textbf{435.9} &  $2.2\times$ \\
$16{\times}16{\times}16{\times}16$       &   58.9 &    6.6 &    24.4 &             31.4 &  $0.5\times$ \\
\bottomrule
\end{tabular}
\end{table}

Table~\ref{tab:bench} reports the throughput of the optimised
\textsc{LQCD-Neuron} Wilson Dslash on an \texttt{inf2.8xlarge},
compared against a CPU baseline running the same reference
\texttt{nn.Module} in FP32. Three observations are worth noting.

First, the unoptimised single-RHS Neuron column (column~3) is only
faster than CPU at the smallest volume; by $V=8{\times}8{\times}8{\times}4$ it is
slower by nearly $2\times$, and at $V=16{\times}16{\times}16{\times}16$ by more
than $8\times$. This confirms that the naive Dslash graph is
dispatch- and shape-limited: the work per call is too small to
amortise the $\sim\!1\,$ms NeuronCore launch cost, and the
underlying $3{\times}3$ / $4{\times}4$ einsums leave the MXU idle.
Second, multi-RHS batching closes the gap and overtakes CPU at all
volumes up to $V=16{\times}8{\times}8{\times}8$ ($2.2\times$ advantage with
Multicore), but falls below CPU at $V=16{\times}16{\times}16{\times}16$
($0.5\times$ with Multicore), where the fused $12{\times}12$ kernel
buffers alone exceed NeuronCore SRAM and spill to HBM. Third, the
two-core \texttt{DataParallel} (Multicore) configuration delivers
$6.4$--$13.0\times$ over CPU from $V=4^4$ through
$V=8{\times}8{\times}8{\times}4$. The breakeven volume --- beyond which the
CPU baseline is faster even with all optimisations --- lies between
$V=16{\times}8{\times}8{\times}8$ and $V=16{\times}16{\times}16{\times}16$; bridging
this gap is the primary motivation for the distributed multi-instance
work outlined in Sec.~\ref{sec:outlook}.

We have additionally verified, on CPU and on a 2-core Inf2 in
\texttt{DataParallel} mode, that the BiCGStab and CG host-loop
solvers converge to the expected residuals on a Gaussian random
$\psi$ with anti-periodic temporal boundaries, and that the
gauge-baked, fused, multi-RHS \NEFF\ is invariant under permutation
of right-hand sides to within FP32 round-off.

\section{Limitations and outlook}
\label{sec:outlook}

\paragraph*{Functionality gap with QUDA.}
The current release covers only Wilson and Clover-Wilson fermions, CG
and BiCGStab, periodic / anti-periodic boundaries, and a small set of
gauge observables (plaquette, Wilson action, Polyakov loop,
clover-definition topological charge). Staggered, twisted-mass and
domain-wall formulations~\cite{Kogut:1974ag,Frezzotti:2000nk,Kaplan:1992bt}
are representable in the same \texttt{nn.Module} style but not yet implemented.
Even-odd preconditioning and a multigrid preconditioner are the most
impactful missing components; both are pure-tensor algorithms and
present no fundamental obstacle to the \texttt{neuronx-cc} pipeline.

\paragraph*{Multi-instance scaling.}
Within a single Inf2/Trn1 instance, communication uses NeuronLink
on-chip and \texttt{torch\_neuronx.DataParallel} as shown above. For
multi-instance domain decomposition with halo exchange, AWS exposes
EFA via \texttt{torch.distributed} with the NCCL-style XLA collective
backend; integrating this with a Schwarz-style domain decomposition
is the natural next step and is what is required to address physically
relevant volumes ($V\!\gtrsim\!32{\times}32{\times}32{\times}32$) that exceed single-chip HBM.

\paragraph*{Static shapes.}
The most awkward limitation is the static-shape requirement: each
$(\mathrm{class},\mathrm{shape},\mathrm{dtype},B)$ tuple compiles to a
distinct \NEFF. For a fixed production volume this is amortised once,
but for shape-sweeping workloads (auto-tuning, mixed-volume gauge
ensembles) it is a real cost. \texttt{neuronx-cc} bucketing and
shape-polymorphic XLA are partial mitigations.

\paragraph*{Precision.}
We currently run BF16 for the matvec with FP32 accumulation in the
host-side CG/BiCGStab loop, which suffices to match the existing
unit-test tolerances. Mixed-precision Krylov methods with iterative
refinement~\cite{Clark:2009wm} would map cleanly onto the NeuronCore
MXU's BF16 inputs / FP32 accumulators and are expected to recover full
double-precision residuals at near-BF16 throughput.

\section{Conclusion}
\label{sec:concl}

We have shown that the AoT-compiled, static-graph, matrix-engine
programming model exposed by AWS Trainium and Inferentia is
expressive enough to host the core kernels of Wilson and
Clover-Wilson Lattice QCD, and that with four straightforward
graph-level optimisations - gauge baking, fused $12{\times}12$
spin--colour kernels, multi-RHS batching, and multi-core
\texttt{DataParallel} sharding - a pure-Python library with no
custom kernel code reaches up to $10.9\times$ over a tuned CPU
baseline on $V=4{\times}4{\times}4{\times}4$ in BF16.

\section*{Acknowledgments}
The author thanks the AWS Neuron SDK team for documentation and
support, and the developers of \PyT, XLA, and \textsc{QUDA} whose
designs informed the architecture of this work.

\begin{thebibliography}{99}

\bibitem{Wilson:1974sk}
K.~G.~Wilson,
\emph{Confinement of quarks},
Phys.\ Rev.\ D \textbf{10}, 2445 (1974).

\bibitem{Clark:2009wm}
M.~A.~Clark, R.~Babich, K.~Barros, R.~C.~Brower, and C.~Rebbi,
\emph{Solving Lattice QCD systems of equations using mixed precision
solvers on GPUs},
Comput.\ Phys.\ Commun.\ \textbf{181}, 1517 (2010),
\href{https://arxiv.org/abs/0911.3191}{arXiv:0911.3191}.

\bibitem{Babich:2011np}
R.~Babich, M.~A.~Clark, B.~Jo\'o, G.~Shi, R.~C.~Brower, and S.~Gottlieb,
\emph{Scaling lattice QCD beyond 100 GPUs},
SC'11 Proceedings (2011),
\href{https://arxiv.org/abs/1109.2935}{arXiv:1109.2935}.

\bibitem{Brower:2018kgy}
R.~C.~Brower, M.~A.~Clark, A.~Strelchenko, and E.~Weinberg,
\emph{Multigrid algorithm for staggered lattice fermions},
Phys.\ Rev.\ D \textbf{97}, 114513 (2018),
\href{https://arxiv.org/abs/1801.07823}{arXiv:1801.07823}.

\bibitem{xla}
The XLA Authors,
\emph{XLA: Optimizing compiler for machine learning},
\url{https://openxla.org/xla} (2024).

\bibitem{aws-trainium}
Amazon Web Services,
\emph{AWS Neuron SDK Documentation},
\url{https://awsdocs-neuron.readthedocs-hosted.com/} (2024).

\bibitem{DeGrandRossi}
T.~A.~DeGrand and D.~Toussaint, eds.,
\emph{From Actions to Answers}, World Scientific (1990);
see also F.~Berruto \emph{et al.}, Phys.\ Rev.\ D \textbf{73}, 014503 (2006).

\bibitem{HestenesStiefel}
M.~R.~Hestenes and E.~Stiefel,
\emph{Methods of conjugate gradients for solving linear systems},
J.\ Res.\ Natl.\ Bur.\ Stand.\ \textbf{49}, 409 (1952).

\bibitem{Kogut:1974ag}
J.~B.~Kogut and L.~Susskind,
\emph{Hamiltonian formulation of Wilson's lattice gauge theories},
Phys.\ Rev.\ D \textbf{11}, 395 (1975).

\bibitem{Frezzotti:2000nk}
R.~Frezzotti, P.~A.~Grassi, S.~Sint, and P.~Weisz,
\emph{Lattice QCD with a chirally twisted mass term},
JHEP \textbf{08}, 058 (2001),
\href{https://arxiv.org/abs/hep-lat/0101001}{hep-lat/0101001}.

\bibitem{Kaplan:1992bt}
D.~B.~Kaplan,
\emph{A method for simulating chiral fermions on the lattice},
Phys.\ Lett.\ B \textbf{288}, 342 (1992),
\href{https://arxiv.org/abs/hep-lat/9206013}{hep-lat/9206013}.

\end{thebibliography}

\end{document}
